\section{Round-nine reconstruction: exact phase labels, microscopic mixtures, and commuting contractions}
\label{sec:r9-d1}

The phase decomposition is now part of the microscopic probability space.  A
minimum of phase rates is obtained by contracting an exact labelled mixture,
not by taking the conjugate of a maximum pressure and not by discarding an
$e^{-c\mu}$ remainder, which can change a large-deviation rate.  Phase-local
complex charts are used only for local response and likelihoods; coexistence
is handled by the label variable and allows Lee--Yang pinching.

\subsection{Exact labelled finite-volume laws}

Let $\mathfrak A$ be a compact finite or finite-dimensional phase-label space.
For every $\varepsilon$ let $\rho_\varepsilon$ be a probability on
$\mathfrak A$ and let
$\mathbb P_{\varepsilon,\alpha}$ be a genuine normalized microscopic law on a
common measurable state space $E_\varepsilon$.  The labelled law is
\[
 \overline{\mathbb P}_\varepsilon(d\alpha,dx)
 =\rho_\varepsilon(d\alpha)
  \mathbb P_{\varepsilon,\alpha}(dx).
\]
The physical unlabelled law is its exact projection.  For Sinai phases the
labels are recurrent components/face phases from A3; for hard spheres they are
the B1/B2 prepared exposed phases and coexistence components.  Restricted
partition functions and boundary conditions defining each law are included in
the platform theorem.

\begin{assumption}[Subexponential phase prior]
\label{ass:r9-d1-prior}
The priors satisfy an LDP at speed $\mu_\varepsilon$ with good rate
$J(\alpha)$, and on every compact set of labels their masses are bounded below
by $e^{-o(\mu_\varepsilon)}$ relative to the appropriate neighbourhoods.
There is no discarded signed or total-variation remainder.
\end{assumption}

\begin{theorem}[Joint labelled platform LDP]
\label{thm:r9-d1-joint}
Suppose the platform theorems give, uniformly on compact label sets, a good
conditional LDP for $X_\varepsilon$ under
$\mathbb P_{\varepsilon,\alpha}$ with rate $I_\alpha(x)$, including boundary
face recovery.  Then $(\alpha,X_\varepsilon)$ under
$\overline{\mathbb P}_\varepsilon$ satisfies the good joint LDP
\[
 \overline I(\alpha,x)=J(\alpha)+I_\alpha(x).
\]
Consequently the exact physical mixture satisfies
\[
 I(x)=\inf_{\alpha\in\mathfrak A}
 \{J(\alpha)+I_\alpha(x)\}.
\]
For a finite zero-cost prior this is $\min_\alpha I_\alpha(x)$ because of
the microscopic mixture theorem, not because
$(\max Q_\alpha)^*=\min Q_\alpha^*$.
\end{theorem}

\begin{proof}
For a closed $F\subset\mathfrak A\times E$, cover its compact
$\alpha$-projection by finitely many uniform conditional-LDP charts and combine
the prior upper bound with the conditional upper bounds.  Exponential
tightness makes the rate good.  For an open neighbourhood of $(\alpha,x)$,
use the prior lower bound and the phase-specific microscopic recovery supplied
by A3 or B2/B1, including its face recovery when $x$ lies at the boundary of
the phase mean set.  This proves the joint lower bound.  Contracting the exact
label coordinate gives the physical rate.  No exponentially small term is
ignored: a phase with prior cost $J$ appears explicitly in the infimum.
\end{proof}

\subsection{Phase-local complex response charts}

For a fixed phase label $\alpha$, let
\[
 Z_{\varepsilon,\alpha}(\theta)
 =\mathbb E_{\varepsilon,\alpha}
 e^{\mu_\varepsilon\langle\theta,X_\varepsilon\rangle}.
\]
On a regular phase chart the A2 or B2 transfer/cluster theorem gives a
zero-free complex neighbourhood of a compact real source set and an analytic
pressure $Q_\alpha$.  No one such neighbourhood is asserted across a
coexistence locus.  The physical mixture partition function is the positive
sum
\[
 Z_\varepsilon(\theta)
 =\int Z_{\varepsilon,\alpha}(\theta)\rho_\varepsilon(d\alpha),
\]
and may have complex zeros at distance $O(\mu_\varepsilon^{-1})$ from
coexistence.

\begin{theorem}[Local derivatives and exact finite-centred likelihood]
\label{thm:r9-d1-likelihood}
On a regular fixed-phase chart all pressure derivatives commute with the
microscopic limit.  For a finite-dimensional local direction $h$, the exact
Radon--Nikodym likelihood within that phase is
\[
 L_{\varepsilon,\alpha}(h)
 =\exp\left\{
 \sqrt{\mu_\varepsilon}
 h\cdot(X_\varepsilon-DQ_{\varepsilon,\alpha}(\theta))
 -\mu_\varepsilon\left[
 Q_{\varepsilon,\alpha}(\theta+h/\sqrt\mu)
 -Q_{\varepsilon,\alpha}(\theta)
 -\frac1{\sqrt\mu}DQ_{\varepsilon,\alpha}(\theta)h
 \right]
 \right\}.
\]
It has expectation one exactly and converges jointly with the centred field to
the Cameron--Martin likelihood.  At coexistence, the phase label and its
Bayesian weight remain in the local experiment.
\end{theorem}

\begin{proof}
Cauchy estimates on the phase-local zero-free chart give convergence of all
finite derivatives.  The displayed expression is the exact exponential tilt
normalized by the finite pressure, hence is a mean-one likelihood.  Taylor
expansion at the exact finite mean and the A2/B3 Gaussian limit give the
Cameron--Martin law.  For several active labels, apply the statement
conditionally and update the finite label weights by their exact likelihoods.
No common analytic logarithm is used.
\end{proof}

\subsection{Faces and projective passage}

Let $\pi_m$ be an increasing family of finite continuous projections.  The
platform theorem provides for every label a full finite-dimensional LDP on the
closed phase-attainable set $C_{m,\alpha}$.  Interior points are recovered by
finite tilts.  Boundary faces are recovered by the actual A3 survivor/renewal
construction or the B2/B1 constrained activity construction, not by merely
sending an abstract normal source to infinity.

\begin{theorem}[Labelled projective LDP and topology upgrade]
\label{thm:r9-d1-projective}
The joint labelled finite projections have compatible good rates
\[
 \overline I_m(\alpha,y)
 =J(\alpha)+I_{m,\alpha}(y).
\]
Dawson--G\"artner yields a projective joint LDP with rate
\[
 \overline I(\alpha,x)=
 \sup_m\overline I_m(\alpha,\pi_mx).
\]
If the platform has exponential containment in a stronger weighted topology
and every finite-rate point has a recovery sequence contained in one compact
weighted sublevel, then the LDP upgrades to that topology.  Contracting the
label after this passage gives the physical rate of
Theorem~\ref{thm:r9-d1-joint}; no interchange of an arbitrary projective
supremum with a phase minimum is used.
\end{theorem}

\begin{proof}
Compatibility follows from contraction of the exact joint finite-dimensional
laws.  Dawson--G\"artner gives the projective theorem.  For a target-open set
and a finite-rate point, choose the platform recovery inside a compact rate
sublevel; on that compact the projective and weighted topologies agree, so a
finite projective neighbourhood lies inside the target-open set.  The recovery
law has exponentially negligible mass outside the compact by the platform
source tilt, providing the strong-topology lower bound.  Upper bounds use
exponential containment.  The label is retained until the final ordinary
contraction.
\end{proof}

\subsection{Thin shells across phases}

Let $C(X_\varepsilon)$ be a finite lattice--continuous constraint.  For every
active phase, A2 or B1 supplies a phase-specific exact two-saddle coefficient
$\mathfrak c_{\varepsilon,\alpha}(H,a)$, uniform on the compact active label
set.

\begin{theorem}[Exact labelled thin-shell conditioning]
\label{thm:r9-d1-shell}
For a bounded source $H$, the conditioned physical pressure is obtained from
the exact finite sum/integral
\[
 \frac{
  \int\rho_\varepsilon(d\alpha)
   e^{\mu_\varepsilon q_{\varepsilon,\alpha}(H,a)}
   \mathfrak c_{\varepsilon,\alpha}(H,a)
 }{
  \int\rho_\varepsilon(d\alpha)
   e^{\mu_\varepsilon q_{\varepsilon,\alpha}(0,a)}
   \mathfrak c_{\varepsilon,\alpha}(0,a)
 }.
\]
Laplace's method on the compact label space gives the minimum constrained
phase cost and retains the normalized coefficient weights among tied phases.
Conditioning, typed contraction, local Gaussian tangent, and nonlinear dynamic
programming commute after the label is included in the state.
\end{theorem}

\begin{proof}
Disintegrate the exact labelled law before conditioning and apply the
phase-specific coefficient theorem.  Uniformity allows Laplace's method over
$\alpha$; tied minimizers retain their prior and determinant weights rather
than being merged into a nonexistent common saddle.  Continuous contractions
and dynamic values may then be applied to the joint labelled state.  Local
Gaussian tangents use Theorem~\ref{thm:r9-d1-likelihood} conditionally on the
selected phase.
\end{proof}

\subsection{The standalone commutation theorem}

\begin{theorem}[Exact phase-labelled commutation principle]
\label{thm:r9-d1-main}
For every platform satisfying the explicitly proved A2/A3 or B1/B2/B3/B4
interfaces, the following diagram commutes on the exact labelled law:
\[
 \text{phase preparation}
 \longrightarrow
 \text{path pressure/LDP}
 \longrightarrow
 \text{conditioning and typed contraction}
 \longrightarrow
 \text{Gaussian likelihood and nonlinear value}.
\]
The physical unlabelled statement is obtained by one final contraction of the
phase label.  The theorem permits phase coexistence and Lee--Yang pinching,
never infers a minimum rate from a maximum pressure alone, and does not discard
exponential remainders.
\end{theorem}

\begin{proof}
The joint LDP and physical rate are
Theorem~\ref{thm:r9-d1-joint}; local likelihood and Gaussian passage are
Theorem~\ref{thm:r9-d1-likelihood}; projective/topology passage is
Theorem~\ref{thm:r9-d1-projective}; and exact conditioning is
Theorem~\ref{thm:r9-d1-shell}.  Each operation is performed before contracting
the exact label, so the ordinary contraction and disintegration identities
prove commutation.
\end{proof}
